%% kapa: exam=1 round=35 year=2024 period=1 form=A subject=real_estate subject_ko=부동산학원론 seq=14 number=94
%% source: Q-Net 감정평가사 기출문제 (cst00309) artlSeq=5213229 file=raw/1008/5213229/1교시 문제지 원본.pdf
%% conversion: kordoc 4.24.51 -> tools/gen1.py (auto); review=proofread-sonnet; flags=none
\begin{question}{94}
감정평가 실무기준상 수익방식에 관한 내용으로 옳은 것은?
\begin{choices}
  \item 직접환원법은 복수기간의 순수익을 적절한 환원율로 환원하여 대상물건의 가액을 산정하는 방법을 말한다.
  \item 수익가액이란 수익분석법에 따라 산정된 가액을 말한다.
  \item 순수익은 대상물건에 귀속하는 적절한 수익으로서 가능총수익에서 운영경비를 공제하여 산정한다.
  \item 직접환원법에서 사용할 환원율은 투자결합법으로 구하는 것을 원칙으로 한다.
  \item 할인현금흐름분석법의 적용에 따른 복귀가액은 보유기간 경과 후 초년도의 순수익을 추정하여 최종환원율로 환원한 후 매도비용을 공제하여 산정한다.
\end{choices}
\end{question}
\answer{5}
